\section{Big Data}\label{sec:big-data}

\subsection{The Big Data and AI Market Landscape}\label{subsec:big-data--the-big-data-and-ai-market-landscape}

The first idea to establish is that \textbf{Big Data is not a single technology}. When people say ``\emph{Big Data}'', they may be referring to databases, distributed storage, analytics platforms, cloud infrastructure, machine learning, visualization tools, or complete applications. In reality, these components form an \textbf{ecosystem}. A useful simplified picture is:
\begin{equation*}
    \text{Data Sources}
    \rightarrow
    \text{Data Infrastructure}
    \rightarrow
    \text{Analytics}
    \rightarrow
    \text{Applications}
\end{equation*}
Different companies and technologies operate at different points of this pipeline.

\highspace
\begin{flushleft}
    \textcolor{Green3}{\faIcon{stream} \textbf{Data sources: where everything starts}}
\end{flushleft}
Before we can have ``Big Data'', something must generate data. Modern systems continuously produce information from many heterogeneous sources like web, smartphones, IoT devices, social networks, transactions, sensors, and more.

\highspace
\begin{examplebox}[: E-Commerce data sources]
    For example, imagine an e-commerce service. It might collect: customer profiles, purchases, searches, clicks, page views, product reviews images, payment transactions, delivery information. Notice immediately that \textbf{these are not all the same kind of data}. Some may be highly structured:
\begin{lstlisting}
order_id = 58321
customer_id = 912
price = 49.99
\end{lstlisting}
    While others may be much less structured:
    \begin{center}
        ``\emph{This product arrived late but works very well.}''
    \end{center}
    Or could be images, video, GPS coordinates, logs, etc. This heterogeneity will later become one of the defining properties of Big Data: \textbf{Variety}.
\end{examplebox}

\highspace
\begin{flushleft}
    \textcolor{Green3}{\faIcon{server} \textbf{The infrastructure layer}}
\end{flushleft}
Once data exists, we need technologies capable of \textbf{storing, transporting and processing it}. This is the infrastructure part of the Big Data market. Conceptually:
\begin{equation*}
    \text{Raw Data}
    \rightarrow
    \text{\hl{Infrastructure}}
    \rightarrow
    \text{usable data}
\end{equation*}
Infrastructure contains things such as: distributed storage; databases; distributed computing systems; data ingestion systems; messaging systems; cloud infrastructure. Later in the notes we will encounter technologies belonging to this general world, such as distributed databases and large-scale processing architectures.

\highspace
The fundamental challenge is that a \textbf{traditional architecture} such as \emph{one database on one machine} may eventually become \textbf{insufficient}. Then the architecture can evolve toward something like:
\begin{equation*}
    \begin{array}{ccc}
        \text{Machine 1} & \text{Machine 2} & \text{Machine 3}\\
        \text{data} & \text{data} & \text{data}
    \end{array}
\end{equation*}
But in the other hand, a \textbf{distributed architecture} introduces new questions about partitioning, replication, consistency, failures, and scalability. These questions will be central later in the notes.

\highspace
\begin{flushleft}
    \textcolor{Green3}{\faIcon{code} \textbf{Open-source platforms}}
\end{flushleft}
Another important part of the Big Data ecosystem is the strong presence of \textbf{open-source platforms}. Historically, much of the Big Data revolution was enabled by systems that made distributed storage and computation available without requiring every company to implement these mechanisms from scratch. Instead of an organization implementing distributed file storage, fault recovery, parallel computation, cluster scheduling, data partitioning from zero, it can \textbf{build on an existing platform}.

\highspace
So we can think of these platforms as \textbf{providing a reusable technological foundation}:
\begin{equation*}
    \text{Distributed infrastructure}
    \rightarrow
    \text{platform}
    \rightarrow
    \text{data applications}
\end{equation*}
A famous historical example that we will study later is the \emph{Hadoop} ecosystem.

\highspace
\begin{flushleft}
    \textcolor{Green3}{\faIcon{chart-line} \textbf{Analytics layer}}
\end{flushleft}
Simply storing huge quantities of data does not create value. Imagine a company has 100 TB of customer data. If nobody can derive useful information from it, those 100 TB are mostly a \textbf{cost}.

\highspace
We therefore need another layer:
\begin{equation*}
    \text{Stored Data}
    \rightarrow
    \text{\hl{Analytics}}
    \rightarrow
    \text{Information / Models / Predictions}
\end{equation*}
Analytics can include things such as: statistics; aggregation; data mining; machine learning; recommendation systems; forecasting; anomaly detection.

\highspace
\begin{examplebox}[: Netflix analytics]
    Suppose Netflix records:
    \begin{itemize}
        \item User A watched movie X
        \item User A paused at 32 minutes
        \item User A finished movie Y
        \item User A abandoned movie Z
    \end{itemize}
    The \textbf{infrastructure problem} is: ``\emph{Can we reliably collect and store billions of these events?}''. The \textbf{analytics problem} is: ``\emph{Can we infer something useful from them?}''. For instance:
    \begin{equation*}
        P(\text{user likes movie }m\mid\text{previous behaviour})
    \end{equation*}
    That distinction between \textbf{managing data} and \textbf{extracting information from data} is important.
\end{examplebox}

\highspace
\begin{flushleft}
    \textcolor{Green3}{\faIcon{desktop} \textbf{Application layer}}
\end{flushleft}
At the top, analytics becomes something the user or organization can actually use. For example:
\begin{equation*}
    \text{behaviour data}
    \rightarrow
    \text{recommendation model}
    \rightarrow
    \text{\hl{``Recommended for you''}}
\end{equation*}
The user does not see the distributed storage system, cluster, database, or machine-learning infrastructure. They see the \textbf{application}. Other examples might include:
\begin{equation*}
    \text{traffic data}
    \rightarrow
    \text{traffic prediction}
    \rightarrow
    \text{route recommendation}
\end{equation*}
Or:
\begin{equation*}
    \text{financial transactions}
    \rightarrow
    \text{fraud detection model}
    \rightarrow
    \text{fraud alert}
\end{equation*}
Or:
\begin{equation*}
    \text{sensor readings}
    \rightarrow
    \text{predictive model}
    \rightarrow
    \text{maintenance warning}.
\end{equation*}
So there is a \textbf{transformation}:
\begin{equation*}
    \text{raw data}
    \rightarrow
    \text{processed data}
    \rightarrow
    \text{insight}
    \rightarrow
    \text{action}
\end{equation*}

\highspace
\begin{flushleft}
    \textcolor{Green3}{\faIcon{layer-group} \textbf{Full data-stack providers}}
\end{flushleft}
Not every company operates at only one layer. \hl{Some vendors provide an almost complete \textbf{data stack}.} Instead of offering only a database, or only analytics, they might offer:
\begin{equation*}
    \text{storage}
    +
    \text{processing}
    +
    \text{analytics}
    +
    \text{ML}
    +
    \text{visualization}
    +
    \text{application services}
\end{equation*}
These are the \textbf{full data stack providers}.

\begin{examplebox}
    Cloud platforms are an intuitive modern example of this idea.

    A single provider may offer services for: object storage, databases, stream processing, distributed computing, data warehouses, ML training, model deployment, and business intelligence.
\end{examplebox}

\highspace
Therefore, the \hl{borders between the individual categories of the market are not strict.}

\highspace
\begin{flushleft}
    \textcolor{Green3}{\faIcon{brain} \textbf{Where does AI fit?}}
\end{flushleft}
The AI market is strongly connected to this Big Data ecosystem because modern AI depends heavily on data infrastructure. A simplified relationship is:
\begin{equation*}
    \text{Data}
    \rightarrow
    \text{Data Engineering}
    \rightarrow
    \text{Machine Learning / AI}
    \rightarrow
    \text{Application}
\end{equation*}
For example, suppose we want to train a recommendation model. The ML algorithm itself might be only one piece. Before training it, somebody must:
\begin{enumerate}
    \item Collect user interactions;
    \item Clean them;
    \item Integrate different sources;
    \item Store them;
    \item Transform them into suitable features;
    \item Feed them efficiently to the model.
\end{enumerate}
Thus \hl{AI requires a data system underneath it.} A powerful model with poor data infrastructure is often practically useless.

\highspace
\textcolor{Green3}{\faIcon{not-equal} \textbf{Big Data and AI are therefore not synonyms.}} This distinction is useful.
\begin{itemize}
    \item \textbf{Big Data} is primarily concerned with questions such as: ``\emph{How do we collect, store, integrate, manage and process huge or complex datasets?}''
    \item \textbf{AI/ML} is more concerned with: ``\emph{How do we learn patterns or decision functions from data?}''
\end{itemize}
There is obviously a \textbf{large overlap} because \hl{AI systems often consume enormous quantities of data.} But they solve different categories of problems. For Systems and Methods for Big and Unstructured Data (SMBUD), our main perspective is the \textbf{data-management side}. For example, before asking ``\emph{which neural network should we use?}'', we are much more likely to ask:
\begin{itemize}
    \item \emph{Where does the data live?}
    \item \emph{How is it represented?}
    \item \emph{How is it partitioned?}
    \item \emph{How do we query it?}
    \item \emph{Can the system scale?}
    \item \emph{What consistency guarantees do we have?}
\end{itemize}

\highspace
\begin{examplebox}[: Spotify recommendation pipeline]
    Imagine Spotify wants to recommend songs. At the bottom we have \textbf{data sources}: play events, skips, likes, searches, playlists, artists, songs, devices, locations. These might generate billions of events. Then:
    \begin{enumerate}
        \item \textbf{Data Sources} (user events, song metadata, etc.) generate data.
        \item \textbf{Infrastructure} collects and stores them: Distributed Storage, Databases, Streaming Systems.
        \item \textbf{Analytics layer} looks for patterns: $P(\text{user }u\text{ likes song }s)$.
        \item \textbf{Application} shows: Discover Weekly, Recommended songs, Similar artists.
    \end{enumerate}
    Therefore:
    \begin{equation*}
        \text{User events}
        \to
        \text{Big Data infrastructure}
        \to
        \text{Analytics/AI}
        \to
        \text{Product feat}
    \end{equation*}
    This is essentially the market landscape represented as a technical pipeline.
\end{examplebox}

\highspace
\begin{takeawaysbox}
    We don't need to memorize the Big Data market as a collection of company logos. What matters is recognizing the \textbf{roles} inside the ecosystem:
    \begin{equation*}
        \underbrace{\text{Data Sources}}_{\text{generate}}
        \rightarrow
        \underbrace{\text{Infrastructure}}_{\text{store/process}}
        \rightarrow
        \underbrace{\text{Analytics}}_{\text{understand/predict}}
        \rightarrow
        \underbrace{\text{Applications}}_{\text{create value}}
    \end{equation*}
    With \textbf{open-source platforms} providing many of the underlying technological building blocks and \textbf{full-stack providers} potentially spanning several or all of these layers.

    \highspace
    And the AI market sits naturally on top of and alongside this ecosystem because \hl{without a scalable data pipeline, scalable AI is extremely difficult}.
\end{takeawaysbox}